\documentclass[10pt,a4paper]{amsart}
\usepackage[margin=19mm]{geometry}
\usepackage{amsmath,amssymb,mathtools}
\usepackage[scr=rsfs,cal=euler]{mathalfa}
\usepackage{xfrac}
\usepackage[unicode,hidelinks,bookmarks=false]{hyperref}
\newcommand{\ie}{i.e.\ }
\newcommand{\cf}{cf.\ }
\newcommand{\resp}{respectively\ }
\newcommand{\Subsection}{\subsection}
\providecommand{\nobreakdash}{\nobreak\hbox{-}}
\setlength{\emergencystretch}{2em}
\multlinegap0pt
\usepackage{algorithm}\usepackage{algpseudocode}
\usepackage[shortlabels]{enumitem}
\usepackage{bm}
\newtheorem{prop}{Proposition}
\def\BR{{\mathbb{R}}}
\title{Decentralized Stochastic Nonconvex Optimization under~the~$(L_0,L_1)$-Smoothness}
\author{Luo Luo, Xue Cui, Tingkai Jia, Cheng Chen}
\date{Source: arXiv:2509.08726v3 (CC BY 4.0)}
\begin{document}
\maketitle

\noindent\emph{Source and licence.} Decentralized Stochastic Nonconvex Optimization under~the~$(L_0,L_1)$-Smoothness. Version arXiv:2509.08726v3 (CC BY 4.0), distributed by arXiv under the Creative Commons Attribution 4.0 International licence, http://creativecommons.org/licenses/by/4.0/, as recorded in the arXiv licence metadata for this exact version and shown on the abstract page https://arxiv.org/abs/2509.08726v3. Full licence notice: \texttt{licenses/arxiv-2509.08726-CC-BY-4.0.txt}.\par\smallskip

\noindent\textit{Editorial scope.} Counted unit: the article's prop prop:example (source0.tex lines 840--846). The article's complete original proof of it is delivered with it (source0.tex lines 848--891, one \texttt{proof} environment of 44 lines). No mathematics is written, repaired or re-derived: the statement and the proof are the article's own lines, in the article's own notation, and no retained line is substituted or re-flowed.  No further source-local context is needed: every cross-reference of every retained line resolves inside the two delivered spans.  Nothing was omitted from any retained span, so the package carries no omission record.  No retained line contains a citation, so the wrapper carries no bibliography and no bibliography entry.  No retained line contains a figure environment, an image-inclusion command or drawing code, so no image is delivered and the package declares no asset dependency.  Editorial formatting only, in addition: an AMS \texttt{amsart} preamble that restates the article's own declarations and macros used by the retained lines, namely the environments \texttt{prop} on separate counters, together with the standard packages those lines need.  The retained environments print in this document as Proposition 1 in order of appearance, whereas the article prints the same items with the numbering of its own full document.  The complete native source of the article as distributed in the arXiv e-print for this exact version is bundled unchanged as \texttt{sources/184843/source0.tex}.
\begin{prop}\label{prop:example}
We consider functions
\begin{align*}
    f_1(x) = \exp(Lx),~~f_2(x) = \exp(-Lx),~~\text{and}~~f(x) = \frac{f_1(x)+f_2(x)}{2}
\end{align*}
for some $L>0$ and let $L_1=L/\log 2$, then both functions $f_1$ and $f_2$ is $(0,L_1)$-smooth while the function $f$ is not $(0,L_1)$-smooth.
\end{prop}

\begin{proof}
We can verify that the derivatives of functions $f_1$, $f_2$, and~$f$ have the form of $f'_1(x) = L\exp(Lx)$, $f'_2(x) = -L\exp(-Lx)$, and $f'(x) = (L\exp(Lx)-L\exp(-Lx))/{2}$.
Note that the function
\begin{align*}
    h(z) = \frac{|1-\exp(Lz)|}{L|z|} 
\end{align*}
is increasing.
Taking $z=y-x\neq 0$, we have
\begin{align}\label{eq:example-h}
\begin{split}    
  \frac{|1-\exp(L(y-x))|}{L|y-x|} 
= h(y-x)  
\leq   h\left(\frac{1}{L_1}\right)  
= \frac{|1-\exp(L/L_1)|}{L/L_1}
\end{split}
\end{align}
for all $x,y\in\BR$ with $|y-x|\in[-1/L_1, 1/L_1]$.

For the function $f_1$, we have
\begin{align*}
  &  |f'_1(x) - f'_1(y)|  
=   |L\exp(Lx) - L\exp(Ly)| \\
= & L\exp(Lx)|1-\exp(L(y-x))|  \\
\leq &  L\exp(Lx) \cdot L_1|1-\exp(L/L_1)||y-x|  \\
= & L \exp(Lx) |y-x| L_1 
=   L_1 f_1'(x)|x-y|,
\end{align*}
where the inequality is based on equation (\ref{eq:example-h}) and the last two steps are based on $f'_1(x) = L\exp(Lx)$ and the setting $L_1=L/\log 2$.
Therefore, we have shown that $f_1$ is $(0,L_1)$-smooth. 
Similarly, we can also show that $f_2$ is $(0,L_1)$-smooth. 

Let $x=0$ and $y=(\log 2)/L$, which holds $|x-y|=1/L_1$.
Note that we also have
\begin{align*}
    |f'(x) - f'(y)|  
=& \frac{L}{2}\left|\exp(Lx)\!-\!\exp(-Lx) -(\exp(Ly)\!-\!\exp(-Ly))\right| \\
=& \frac{L}{2}\left|\exp(0)-\exp(0) -\left(2-\frac{1}{2}\right)\right| = \frac{3L}{4}
\end{align*}
and
\begin{align*}
    L_1f'(x)|y-x| = \frac{L_1L(\exp(Lx)-\exp(-Lx))|y-x|}{2} = 0.
\end{align*}
Therefore, the inequality $|f'(x) - f'(y)| \leq L_1 f'(x)|y-x|$ does not hold, which means the function~$f$ is not $(0,L_1)$-smooth.
\end{proof}

\end{document}
